\documentclass[12pt]{article}
\usepackage[utf8]{inputenc}
\usepackage[T1]{fontenc}
\usepackage[margin=1in]{geometry}
\usepackage{setspace}
\usepackage{xcolor}
\usepackage[colorlinks=true, urlcolor=blue!60!black, linkcolor=blue!60!black]{hyperref}
\setlength{\parindent}{2em}
\setlength{\parskip}{1.25em}
\renewcommand{\baselinestretch}{1.0}
\title{\textit{CFPB v. Community Financial Services Ass'n of America, Ltd.}, 601 U.S. 416 (2024)}
\author{Case Brief}
\date{}

\begin{document}
\maketitle

\section{Case Metadata}
\begin{description}
  \item[Citation:] 601~U.S.~416 (2024) (No.~22--448).
  \item[Argued / Decided:] Argued October~3, 2023; Decided May~16, 2024.
  \item[Vote:] 7--2 (reversing and remanding 51~F.4th~616 (5th~Cir.~2022)).
  \item[Author:] Justice Thomas (joined by Chief Justice Roberts and Justices Sotomayor, Kagan, Kavanaugh, Barrett, and Jackson).
  \item[Concurrences / Dissents:] Justice Kagan filed a concurring opinion (joined by Justices Sotomayor, Kavanaugh, and Barrett); Justice Jackson filed a concurring opinion; Justice Alito filed a dissenting opinion (joined by Justice Gorsuch).
  \item[Chapter / Topic:] Separation of Powers and the Executive --- Appropriations Clause (Art.~I, \S~9, cl.~7); Power of the Purse; Independent Agency Funding.
  \item[Primary Source URLs:] \href{https://www.supremecourt.gov/opinions/23pdf/22-448_o7jp.pdf}{SupremeCourt.gov Slip Opinion}; \href{https://supreme.justia.com/cases/federal/us/601/22-448/}{Justia U.S. Supreme Court Center}.
\end{description}

\section{Facts and Procedural History}
In response to the 2008 financial crisis, Congress enacted the Dodd-Frank Wall Street Reform and Consumer Protection Act of 2010, 124~Stat.~1376, establishing the Consumer Financial Protection Bureau (CFPB) as an independent regulator within the Federal Reserve System, 12~U.S.C.~\S~5491(a). Congress consolidated in the CFPB the administration of 18 existing consumer financial protection statutes (including the Fair Debt Collection Practices Act, the Fair Credit Reporting Act, and the Home Mortgage Disclosure Act of 1975) and vested the agency with rulemaking, enforcement, and adjudicatory authority over ``unfair, deceptive, or abusive'' acts or practices, \S\S~5512(a), 5531, 5536(a)(1)(B).

Unlike most federal agencies funded through annual congressional appropriations, Congress provided the CFPB a standing funding mechanism outside the annual appropriations process. Under 12~U.S.C.~\S\S~5497(a)(1)--(2), the CFPB Director may requisition each year from ``the combined earnings of the Federal Reserve System'' the amount the Director determines ``to be reasonably necessary to carry out'' the Bureau's duties, subject to an inflation-adjusted cap of 12\% of the Federal Reserve System's total FY~2009 operating expenses (approximately \$734~million in FY~2022). The CFPB may also retain and invest unspent funds from year to year without returning surpluses to the Treasury, \S\S~5497(b)--(c). (In \textit{Seila Law LLC v. CFPB}, 591~U.S.~197 (2020), the Supreme Court had previously severed the CFPB Director's for-cause removal protection under Article~II, leaving the funding mechanism untouched.)

In 2017, the CFPB promulgated its Payday Lending Rule, 12~C.F.R.~pt.~1041 (2018), which restricted lenders from attempting to withdraw loan payments via preauthorized account access after two consecutive failed attempts. Two trade associations representing payday lenders and credit-access businesses---the Community Financial Services Association of America, Ltd. (CFSA) and the Consumer Service Alliance of Texas---sued in the U.S. District Court for the Western District of Texas, challenging the rule on statutory and constitutional grounds, including that the CFPB's funding structure violated the Appropriations Clause.

The district court granted summary judgment to the CFPB, concluding that because a statute authorized disbursements up to a cap, ``there is no Appropriations Clause issue.'' 558~F.Supp.3d~350, 364 (W.D.~Tex.~2021). The Fifth Circuit reversed and vacated the Payday Lending Rule, 51~F.4th~616 (5th~Cir.~2022), holding that the CFPB's ``self-actualizing, perpetual funding mechanism'' violated the Appropriations Clause and constitutional separation of powers by ceding Congress's affirmative duty to maintain periodic fiscal control. The Supreme Court granted certiorari on February~27, 2023.

\section{Issue(s) Presented}
Whether Congress's statutory authorization allowing the CFPB to draw funds from the combined earnings of the Federal Reserve System, in an amount determined by the Director up to an inflation-adjusted cap and without annual congressional renewal, complies with the Appropriations Clause of Article~I, Section~9, Clause~7.

\section{Holding and Rule}
\textbf{Holding:} Yes. Congress's statutory funding mechanism for the CFPB satisfies the Appropriations Clause.

\textbf{Rule:} Under the Appropriations Clause, an ``Appropriatio[n] made by Law'' requires only a statute that identifies a specified source of public funds and authorizes the expenditure of those funds for designated purposes. The Clause does not require annual renewal, line-item specificity, or a fixed dollar sum rather than a statutory cap.

\section{Reasoning}
Writing for a seven-Justice majority, Justice Thomas grounded the Court's analysis in constitutional text, pre-founding history, and early congressional practice:
\begin{enumerate}
  \item \textbf{Threshold Applicability:} Because surplus earnings of the Federal Reserve System would otherwise be deposited into the general fund of the Treasury (12~U.S.C.~\S~289(a)(3)(B)), the CFPB's funding constitutes ``Money \ldots\ drawn from the Treasury'' and must comply with the Appropriations Clause.
  \item \textbf{Textual Meaning of ``Appropriation'':} Consulting constitutional context and Founding-era dictionaries (Webster, Johnson, Ash, Dyche \& Pardon), the Court concluded that an ``appropriation'' meant assigning or setting apart particular public funds for designated uses. Notably, Article~I, \S~8, cl.~12 imposes a two-year durational limit on army appropriations (``no Appropriation of Money to that Use shall be for a longer Term than two Years''), demonstrating by negative implication that the Constitution does not require time-limited appropriations for civilian agencies.
  \item \textbf{Pre-Founding English and Colonial Practice:} While the struggle culminating in the Glorious Revolution established parliamentary supremacy over the purse, Parliament did not require every appropriation to be annual or line-itemized. It enacted standing life grants (such as the civil list) and permissive ``any Sum not exceeding'' caps granting the Crown discretion to spend less than the maximum. Colonial and early state legislatures similarly enacted open-ended, discretionary appropriations.
  \item \textbf{Practice of the First Congress:} Early federal appropriations confirmed wide legislative flexibility. The First Congress enacted lump-sum ``sums not exceeding'' grants (giving the Executive discretion over how much to spend up to a cap) and established standing, indefinite, fee- and commission-based funding mechanisms for the Customs Service (Act of July~31, 1789) and the Post Office (Act of Feb.~20, 1792) that operated without annual appropriations.
  \item \textbf{Rejection of Respondents' Three Objections:} First, the CFPB does not unilaterally set its own budget; Congress set the ceiling via an inflation-adjusted cap, and executive discretion to draw \textit{less} than a cap mirrors early ``sums not exceeding'' laws. Second, the Appropriations Clause imposes no durational limit outside the Army Clause. Third, the Appropriations Clause is phrased as a negative limitation in Article~I, \S~9 rather than a freestanding structural mandate requiring annual review; other constitutional provisions and political checks police broader separation-of-powers concerns.
\end{enumerate}

\section{Concurring and Dissenting Opinions}
\begin{itemize}
  \item \textbf{Justice Kagan, concurring (joined by Justices Sotomayor, Kavanaugh, and Barrett):} Joined the majority in full, but emphasized post-ratification ``long settled and established practice'' across the 19th, 20th, and 21st centuries. She noted that Congress has consistently employed lump-sum grants, standing appropriations (accounting for nearly two-thirds of the FY~2022 federal budget), and assessment-based funding for every federal bank regulator (including the OCC and Federal Reserve Board). The CFPB's mechanism ``would have fit right in'' at the Founding and at every point since.
  \item \textbf{Justice Jackson, concurring:} Agreed that the plain text of the Appropriations Clause resolves the case and urged judicial restraint: when the Constitution's text imposes no limit on a coordinate branch's choice of institutional design, Article~III courts should not invent unstated constraints to second-guess Congress's response to the 2008 financial crisis.
  \item \textbf{Justice Alito, dissenting (joined by Justice Gorsuch):} Argued that ``Appropriations'' is a historically freighted term of art rooted in centuries of parliamentary resistance to Stuart monarchs and colonial assemblies' control over royal governors. In his view, the CFPB's unprecedented combination of six insulating features---perpetuity, self-selected draw amounts, indirect funding from federally chartered Federal Reserve Banks, retention and investment of surpluses, exemption from Appropriations Committee review, and vast regulatory power---``turns the Appropriations Clause into a minor vestige.'' He added that \textit{Seila Law} worsened the problem by placing this fiscally autonomous agency under direct presidential removal control.
\end{itemize}

\section{Key Quotations}
\begin{quote}
``Under the Appropriations Clause, an appropriation is simply a law that authorizes expenditures from a specified source of public money for designated purposes.'' \\
--- Justice Thomas, \href{https://supreme.justia.com/cases/federal/us/601/22-448/}{601~U.S.~at~424} (Majority Opinion).
\end{quote}

\begin{quote}
``Based on the Constitution's text, the history against which that text was enacted, and congressional practice immediately following ratification, we conclude that appropriations need only identify a source of public funds and authorize the expenditure of those funds for designated purposes to satisfy the Appropriations Clause.'' \\
--- Justice Thomas, \href{https://supreme.justia.com/cases/federal/us/601/22-448/}{601~U.S.~at~426} (Majority Opinion).
\end{quote}

\begin{quote}
``Whether or not the CFPB's mechanism has an exact replica, its essentials are nothing new. And it was devised more than two centuries into an unbroken congressional practice, beginning at the beginning, of innovation and adaptation in appropriating funds.'' \\
--- Justice Kagan, \href{https://supreme.justia.com/cases/federal/us/601/22-448/}{601~U.S.~at~446} (Concurring Opinion).
\end{quote}

\begin{quote}
``When the Constitution's text does not provide a limit to a coordinate branch's power, we should not lightly assume that Article~III implicitly directs the Judiciary to find one.'' \\
--- Justice Jackson, \href{https://supreme.justia.com/cases/federal/us/601/22-448/}{601~U.S.~at~447} (Concurring Opinion).
\end{quote}

\begin{quote}
``Unfortunately, today's decision turns the Appropriations Clause into a minor vestige. The Court upholds a novel statutory scheme under which the powerful Consumer Financial Protection Bureau (CFPB) may bankroll its own agenda without any congressional control or oversight. \ldots\ It is not an exaggeration to say that the CFPB enjoys a degree of financial autonomy that a Stuart king would envy.'' \\
--- Justice Alito, \href{https://supreme.justia.com/cases/federal/us/601/22-448/}{601~U.S.~at~448, 467} (Dissenting Opinion).
\end{quote}

\section{Doctrinal Significance and Modern Legacy (Through 2026)}
\textit{CFPB v. CFSA} averted a potential fiscal earthquake by insulating permanent and self-funded federal programs---including Social Security, Medicare, the Federal Reserve Board, the OCC, the FDIC, and the FHFA---from Appropriations Clause invalidation. Methodologically, the 7--2 coalition (with Justice Thomas writing for a cross-ideological majority and Justices Kagan, Kavanaugh, and Barrett concurring on historical liquidation) illustrated that textualist originalism does not invariably produce anti-administrative outcomes when Founding-era statutory practice demonstrates broad legislative flexibility.

Through 2026, the decision's legacy has unfolded across three major fronts:
\begin{itemize}
  \item \textbf{Nondelegation Spillover:} In \href{https://www.law.cornell.edu/supremecourt/text/24-354}{\textit{FCC v. Consumers' Research}, 606~U.S.~\_\_\_ (June~27, 2025)}, the Supreme Court cited \textit{CFPB v. CFSA} (601~U.S.~at~446) while reversing the Fifth Circuit and upholding the FCC's Universal Service Fund contribution scheme against public and private nondelegation challenges.
  \item \textbf{The 2025--2026 ``Combined Earnings'' Statutory Battle:} Because the Supreme Court upheld only the constitutional validity of \S~5497, litigation pivoted to the statute's text. After the Federal Reserve System began operating at a net loss in 2022, a November~2025 DOJ Office of Legal Counsel opinion concluded that ``combined earnings'' in 12~U.S.C.~\S~5497(a)(1) means net profits rather than gross revenues, prompting Acting Director Russell Vought to assert that the CFPB could not lawfully draw funds and warning of a funding lapse (\href{https://daltontomich.com/yet-another-way-to-kill-the-cfpb/}{Dalton \& Tomich analysis, Nov.~2025}). Federal district courts rejected that interpretation and ordered the Bureau to request Federal Reserve funding, prompting a CFPB appeal to the Ninth Circuit in May~2026 (\href{https://www.consumerfinancemonitor.com/2026/05/12/cfpb-surprisingly-appeals-judge-davilas-funding-decision-to-the-ninth-circuit/}{Consumer Finance Monitor, May~2026}).
  \item \textbf{Executive Impoundment Debates:} During 2025--2026 disputes over whether Article~II authorizes the President to impound congressionally appropriated funds notwithstanding the Impoundment Control Act of 1974, scholars and litigants have invoked \textit{CFPB v. CFSA}'s opening declaration that ``[o]ur Constitution gives Congress control over the public fisc'' as a reaffirmed baseline of legislative fiscal supremacy (\href{https://en.wikipedia.org/wiki/Unitary_executive_theory}{Unitary Executive Theory overview}).
\end{itemize}

\section{Pedagogical Critiques and Exam Hypotheticals}
\textbf{Critiques:} Critics of the majority (echoing Justice Alito) argue that the ``source plus purpose'' test is formalistic and nearly toothless: it appears to allow Congress to enact a perpetual, uncapped appropriation authorizing the President to spend ``any sums necessary'' from the general Treasury for any lawful executive purpose, gutting the functional check of the purse. Conversely, defenders note that any future Congress can repeal or amend a standing appropriation by ordinary legislation (subject to the veto), and that Justice Alito's multi-factor ``novelty'' test lacks a textual anchor in Article~I, \S~9, cl.~7 and would embroil courts in arbitrary line-drawing over how much fiscal independence is ``too much.''

\textbf{Exam Hypothetical:} Suppose Congress enacts the ``National AI Safety Authority Act,'' authorizing a single-director agency to draw ``whatever amounts the Director deems necessary, without any dollar or percentage cap, from the general Treasury of the United States in perpetuity, to promote safe artificial intelligence.'' Regulated AI firms challenge the funding mechanism under the Appropriations Clause.
\begin{itemize}
  \item \textbf{Government's Argument (Under \textit{CFPB v. CFSA}):} The statute satisfies Justice Thomas's two-part test: it identifies a source (the general Treasury) and a designated purpose (promoting safe AI). Neither durational limits nor numerical caps are textually required outside the Army Clause, and Congress has long enacted standing ``sums necessary'' mandatory appropriations (as Justice Kagan highlighted).
  \item \textbf{Opposite Side Contrast (Distinguishing \textit{CFPB v. CFSA}):} Challengers can exploit the majority's repeated emphasis on the CFPB's \textit{statutory cap} (``Congress determined the amount of the Bureau's annual funding by imposing a statutory cap,'' 601~U.S.~at~438) and its analogy to early ``sums not exceeding'' grants. Unlike the CFPB or the fee-funded Post Office/Customs Service, an \textit{uncapped} draw from the \textit{general Treasury} paired with an open-ended statutory mandate (``promote safe AI'') combines an Appropriations Clause defect with a severe Article~I nondelegation problem---a separate constitutional constraint that Justice Thomas explicitly left open (601~U.S.~at~440--441).
\end{itemize}

\end{document}
